% !TEX root = paper-node2vec.tex

We formulate feature learning in networks as a maximum likelihood optimization problem. Let $G = (V, E)$ be a given network. Our analysis is general and applies to any (un)directed, (un)weighted network. 
Let $f: V \rightarrow \mathbb{R}^d$ be the mapping function from nodes to feature representaions we aim to learn for a downstream prediction task. Here $d$ is a parameter specifying the number of dimensions of our feature representation. % and is fixed for all $u \in V$. 
Equivalently, $f$ is a matrix of size $|V| \times d$ parameters. 
%
For every source node $u \in V$, we define $N_S(u) \subset V$ as a {\em network neighborhood} of node $u$ generated through a neighborhood sampling strategy $S$. 

We proceed by extending the Skip-gram architecture to networks \cite{word2vec,deepwalk}. 
We seek to optimize the following objective function, which maximizes the log-probability of observing a network neighborhood $N_S(u)$ for a node $u$ conditioned on its feature representation, given by $f$:
%
\vspace{-0.05in}
\begin{equation}\label{eq:rawObjFn}
\begin{aligned}
& \underset{f}{\text{max}}
& &  \sum_{u \in V}\log Pr(N_S(u)| f(u)).
\end{aligned}
\end{equation}

In order to make the optimization problem tractable, we make two standard assumptions:
\begin{itemize}[noitemsep,nolistsep]
	%\denselist
	\item Conditional independence. We factorize the likelihood by assuming that the likelihood of observing a neighborhood node is independent of observing any other neighborhood node given the feature representation of the source:
	\begin{align}
		Pr(N_S(u)| f(u)) = \prod_{n_i \in N_S(u)} Pr(n_i | f(u)). \nonumber
	\end{align}
	\item Symmetry in feature space. A source node and neighborhood node have a symmetric effect over each other in feature space. Accordingly, we model the conditional likelihood of every source-neighborhood node pair as a softmax unit parametrized by a dot product of their features:

	\begin{align}
	Pr(n_i | f(u)) = \frac{\exp(f(n_i)\cdot f(u))}{\sum_{v \in V} \exp(f(v)\cdot f(u))}. \nonumber
	\end{align}
\end{itemize}
%
With the above assumptions, the objective in Eq.~\ref{eq:rawObjFn} simplifies to:
%
\vspace{-0.05in}
\begin{equation}\label{eq:finalObjFn}
	\begin{aligned}
	& \underset{f}{\text{max}}
	& &    \sum_{u \in V} \bigg[- \log Z_u + \sum_{n_i \in N_S(u)} f(n_i)\cdot f(u)\bigg].
	\end{aligned}
\end{equation}
%
The per-node partition function, $Z_u=\sum_{v \in V} \exp(f(u)\cdot f(v))$, is expensive to compute for large networks and we approximate it using negative sampling~\cite{word2vec2}. We optimize Eq.~\ref{eq:finalObjFn} using stochastic gradient ascent over the model parameters defining the features $f$.

Feature learning methods based on the Skip-gram architecture have been originally developed in the context of natural language~\cite{word2vec}. Given the linear nature of text, the notion of a neighborhood can be naturally defined using a sliding window over consecutive words. Networks, however, are not linear, and thus a richer notion of a neighborhood is needed.
%
To resolve this issue, we propose a randomized procedure that samples many different neighborhoods of a given source node $u$. The neighborhoods $N_S(u)$ are not restricted to just immediate neighbors but can have vastly different structures depending on the sampling strategy $S$. 
%At a high level, our sampling strategies are based on the intuition that the context of a node is best represented by its nearby nodes~\cite{harris-1954}. Hence, we will define our context set $N_S(u)$ of node $u$ to be a neighborhood of nodes around $u$.

%\jure{Either use context set or neighborhood set but not both. I kind-of like context set better.}

%Unlike NLP tasks where we have access to a document corpus, there is no empirical distribution over nodes in networks which can be used to estimate $f$. Equivalently, we are not explicitly provided context nodes for any source node. To resolve this issue, we will sample context nodes for every source node from the underlying network. At a high level, the sampling strategies of interest are all based on the intuitive distributional hypothesis analogy \cite{harris-1954} which posits that the context of a node is best represented by nearby nodes. Hence, we will define our context set $N_S(u)$ to be a neighborhood of nodes around the source. Our neighborhoods are not strictly restricted to just immediate neighbors, and can have vastly different interpretations based on our sampling strategy $S$.

\subsection{Classic search strategies}
\label{sec:bfsdfs}

% \begin{figure}[t]
% 	\centering
% 	\includegraphics[width=0.45\textwidth]{FIG/bfsdfs}
% 	\caption{BFS and DFS search strategies from node $u$ ($k=3$).
% 	%\jure{Update for $k=3$. Make the graph prettier, remove some nodes. Add legend: BFS: arrow type, DFS: arrow type. Also, could we make the graph such that we could illustrate structural vs. homiphilic equvivalence. Basically have 2 communities, each with a central node.}
% 	}
% 	\label{fig:bfsdfs}
% \end{figure}

%Building on the above intuition we aim to represent each node with a set of its neighboring nodes. Our goal is to define a flexible network neighborhood sampling strategy that will allow us to generate rich embeddings and capture various types of network equivalences. Importantly, to be able to fairly compare different sampling strategies $S$, we shall constrain the size of the neighborhood set $N_S$ to $k$ nodes and then sample multiple sets for a single node $u$.

We view the problem of sampling neighborhoods of a source node as a form of local search. Figure \ref{fig:bfsdfs} shows a graph, where given a source node $u$ we aim to generate (sample) its neighborhood $N_S(u)$. Importantly, to be able to fairly compare different sampling strategies $S$, we shall constrain the size of the neighborhood set $N_S$ to $k$ nodes and then sample multiple sets for a single node $u$. Generally, there are two extreme sampling strategies for generating neighborhood set(s) $N_S$ of $k$ nodes:
\begin{itemize}[noitemsep,nolistsep]
	\item \textbf {Breadth-first Sampling (BFS)} The neighborhood $N_S$ is restricted to nodes which are immediate neighbors of the source. For example, in Figure~\ref{fig:bfsdfs} for a neighborhood of size $k=3$, BFS samples nodes $s_1$, $s_2$, $s_{3}$.
	\item \textbf {Depth-first Sampling (DFS)} The neighborhood consists of nodes sequentially sampled at increasing distances from the source node. In Figure~\ref{fig:bfsdfs}, DFS samples $s_{4}$, $s_{5}$, $s_{6}$. 
\end{itemize}

%\jure{Move the computational issues somehwere else -- they distract here}

%%%%%%%% JURE

The breadth-first and depth-first sampling represent extreme scenarios in terms of the search space they explore leading to interesting implications on the learned representations. 

In particular, prediction tasks on nodes in networks often shuttle between two kinds of similarities: homophily and structural equivalence~\cite{latentfactor}. Under the homophily hypothesis~\cite{fortunato09community, yang14ieee} nodes that are highly interconnected and belong to similar network clusters or communities should be embedded closely together (\eg, nodes $s_1$ and $u$ in Figure~\ref{fig:bfsdfs} belong to the same network community). In contrast, under the structural equivalence hypothesis~\cite{henderson-kdd2012} nodes that have similar structural roles in networks should be embedded closely together (\eg, nodes $u$ and $s_6$ in Figure~\ref{fig:bfsdfs} act as hubs of their corresponding communities). Importantly, unlike homophily, structural equivalence does not emphasize connectivity; nodes could be far apart in the network and still have the same structural role.
In real-world, these equivalence notions are not exclusive; networks commonly exhibit both behaviors where some nodes exhibit homophily while others reflect structural equivalence.
%Also, it is possible that the roles in structural equivalence could overlap.

We observe that BFS and DFS strategies play a key role in producing representations that reflect either of the above equivalences. In particular, the neighborhoods sampled by BFS lead to embeddings that correspond closely to structural equivalence. Intuitively, we note that in order to ascertain structural equivalence, it is often sufficient to characterize the local neighborhoods accurately. For example, structural equivalence based on network roles such as bridges and hubs can be inferred just by observing the immediate neighborhoods of each node. By restricting search to nearby nodes, BFS achieves this characterization and obtains a microscopic view of the neighborhood of every node. Additionally, in BFS, nodes in the sampled neighborhoods tend to repeat many times. This is also important as it reduces the variance in characterizing the distribution of 1-hop nodes with respect the source node. However, a very small portion of the graph is explored for any given $k$.

The opposite is true for DFS which can explore larger parts of the network as it can move further away from the source node $u$ (with sample size $k$ being fixed). In DFS, the sampled nodes more accurately reflect a macro-view of the neighborhood which is essential in inferring communities based on homophily. However, the issue with DFS is that it is important to not only infer which node-to-node dependencies exist in a network, but also to characterize the exact nature of these dependencies. This is hard given we have a constrain on the sample size and a large neighborhood to explore, resulting in high variance. Secondly, moving to much greater depths leads to complex dependencies since a sampled node may be far from the source and potentially less representative. 

\subsection{node2vec}

Building on the above observations, we design a flexible neighborhood sampling strategy which allows us to smoothly interpolate between BFS and DFS. We achieve this by developing a flexible biased random walk procedure that can explore neighborhoods in a BFS as well as DFS fashion.

\subsubsection{Random Walks}
Formally, given a source node $u$, we simulate a random walk of fixed length $l$. Let $c_i$ denote the $i${th} node in the walk, starting with $c_0 = u$. Nodes $c_i$ are generated by the following distribution:
%
\[
	P(c_i = x \mid c_{i-1} = v) =
	\begin{cases}
	\frac{\pi_{vx}}{Z} & \text{if } (v,x) \in E \\
	0 & \text{otherwise}
	\end{cases}
\]
%
where $\pi_{vx}$ is the unnormalized transition probability between nodes $v$ and $x$, and $Z$ is the normalizing constant. 
%The simplest way would be to transition based on the weights of the edges in the graph ($\pi_{vx} = w_{vx}$). (In case of unweighted graphs $w_{vx}=1$.) However, we want to adaptively change transition probabilities $\pi_{vx}$ based on the network structure and this way we can guide the random walk to explore different types of network neighborhoods.

\subsubsection{Search bias \large{$\alpha$}}
%A naive way to bias our random walks would  be to sample the next node based on the static edge weight $w_{vx}$. (Note that for unweighed graphs all $w_{vx}=1$.) However, this does not allow us to account for the network structure and guide our search procedure to explore different types of network neighborhoods. 
The simplest way to bias our random walks would be to sample the next node based on the static edge weights $w_{vx}$ \ie, $\pi_{vx} = w_{vx}$. (In case of unweighted graphs $w_{vx}=1$.) However, this does not allow us to account for the network structure and guide our search procedure to explore different types of network neighborhoods. 
Additionally, unlike BFS and DFS which are extreme sampling paradigms suited for structural equivalence and homophily respectively, our random walks should accommodate for the fact that these notions of equivalence are not competing or exclusive, and real-world networks commonly exhibit a mixture of both. 

We define a 2$^{\textrm{nd}}$ order random walk with two parameters $p$ and $q$ which guide the walk: Consider a random walk that just traversed edge $(t,v)$ and now resides at node $v$ (Figure~\ref{fig:walk}). The walk now needs to decide on the next step so it evaluates the transition probabilities $\pi_{vx}$ on edges $(v,x)$ leading from $v$. We set the unnormalized transition probability to $\pi_{vx} = \alpha_{pq}(t,x)\cdot w_{vx}$, where 
\[
	\alpha_{pq}(t,x) = 
	\begin{cases}
	\frac{1}{p}  & \text{if } d_{tx} = 0\\
	1 & \text{if } d_{tx} = 1\\
	\frac{1}{q} & \text{if } d_{tx} = 2
	\end{cases}
\]
%
and $d_{tx}$ denotes the shortest path distance between nodes $t$ and $x$. Note that $d_{tx}$ must be one of $\{0,1,2\}$, and hence, the two parameters are necessary and sufficient to guide the walk.

\begin{figure}[t]
	\centering
	\includegraphics[width=0.3\textwidth]{FIG/walk}
	\vspace{-0.2cm}
	\caption{Illustration of the random walk procedure in \nodevec. The walk just transitioned from $t$ to $v$ and
	is now evaluating its next step out of node $v$. Edge labels indicate search biases $\alpha$.}
	\label{fig:walk}
	\vspace{-0.4cm}
\end{figure}


Intuitively, parameters $p$ and $q$ control how fast the walk explores and leaves the neighborhood of starting node $u$. In particular, the parameters allow our search procedure to (approximately) interpolate between BFS and DFS and thereby reflect an affinity for different notions of node equivalences.

%We quantify our search bias by introducing a real positive parameter $\alpha$, which in turn is governed by the underlying network $G$ and two hyperparameters, $p$ and $q$. In the random walk setting, this is equivalent to modifying our edge weights through a multiplicative factor $\alpha$. 
%If $u$ is the only node in the walk, then $\alpha=1$. Else, we set $\alpha = g(d_{tv})$ where $t$ is defined as the node preceding $u$ in the random walk and $d_{tv}$ is the distance between $t$ and $v$. 
%, we parametrize $\alpha$ using two hyperparameters,  i.e.
%The hyperparameters $p$ and $q$ (along with standard random walk parameters such as walk length and number of walks per node) allow our search to (approximately) interpolate between BFS and DFS, thereby reflecting an affinity for different notions of equivalences, described below.

\xhdr{Return parameter, $p$} Parameter $p$ controls the likelihood of immediately revisiting a node in the walk. Setting it to a high value ($>\max(q, 1)$) ensures that we are less likely to sample an already-visited node in the following two steps (unless the next node in the walk had no other neighbor). This strategy encourages moderate exploration and avoids $2$-hop redundancy in sampling. On the other hand, if $p$ is low ($<\min(q, 1)$), it would lead the walk to backtrack a step (Figure~\ref{fig:walk}) and this would keep the walk ``local'' close to the starting node $u$.

%To see why, let us assume in Figure \ref{fig:bfsdfs}, we wish to sample nodes directly connected to $u$ which is the starting node in our walk. Without loss of generality, assume that we sample $s_2$. To sample another neighbor of $u$, we could return to $u$ (with high probability) by setting $p$ to a very low value. However, as a result of low $p$, we are likely to sample back $s_2$ from $u$ and consequently go back and forth $u$ and $s_2$ with high probability.

\xhdr{In-out parameter, $q$} Parameter $q$ allows the search to differentiate between ``inward'' and ``outward'' nodes.  Going back to Figure~\ref{fig:walk}, if $q>1$, the random walk is biased towards nodes close to node $t$. Such walks obtain a local view of the underlying graph with respect to the start node in the walk and approximate BFS behavior in the sense that our samples comprise of nodes within a small locality.  
%Although in order to benefit from the $O(\frac{1}{\lfloor\frac{l}{k}\rfloor + 1})$ sampling time of random walks, we do not reduce the walk length to diminishingly small values to exactly match BFS. 

In contrast, if $q<1$, the walk is more inclined to visit nodes which are further away from the node $t$. Such behavior is reflective of DFS which encourages outward exploration. However, an essential difference here is that we achieve DFS-like exploration within the random walk framework. Hence, the sampled nodes are not at strictly increasing distances from a given source node $u$, but in turn, we benefit from tractable preprocessing and superior sampling efficiency of random walks. Note that by setting $\pi_{v,x}$ to be a function of the preceeding node in the walk $t$, the random walks are 2$^{\textrm{nd}}$ order Markovian.

%To exploit the full potential of \nodevec for a given domain network, we need to learn $q$ which can be done quickly in a semi-supervised way with little training data since it is essentially a proxy for inferring the global mixture of equivalences in the underlying graph. Since $q$ has a distinct interpretation, it can even be set directly based on domain-expertize if available. As far as the return hyperparameter is concerned, \nodevec is less sensitive and setting $p$ to a moderately high value works well in practice.

% \subsubsection{Node importance}
% \critique{Since PageRank is proportion to degree, this paragraph makes little sense even though it might be working in practice. Maybe we should delete this altogether. Instead, let's move the Les Miserables Case Study to immediately follow the pseudocode in the next subsection}
% Random walks introduce a global bias. In particular, nodes with high degree are more likely to be visited by a random walk and thus will more often appear in the sampled node neighborhoods. The degree bias is not necessarily a bane since nodes with high degree are often the most important nodes in the network. 
% At the same time, there are numerous instances where nodes of low degrees, such as bridge nodes, play important roles and thus one would aim to sample such nodes as well. Thus, it is important to find the right tradeoff between between local and global importance of nodes.

% In our experiments, we find that correcting for degree bias and factoring for PageRanks of nodes leads to better quality embeddings. However, naively incorporating PageRanks of nodes is misleading for learning feature representations. Rather, we normalize the PageRanks by the degree of the node to measure its importance. The importance of nodes defined in terms of normalized PageRanks can be computed efficiently offline and does not contribute to the runtime complexity. 

\xhdr{Benefits of random walks} There are several benefits of random walks over pure BFS/DFS approaches. Random walks are computationally efficient in terms of both space and time requirements. The space complexity to store the immediate neighbors of every node in the graph is $O(|E|)$. 
For 2$^{\textrm{nd}}$ order random walks, it is helpful to store the interconnections between the neighbors of every node, which incurs a space complexity of $O(a^2|V|)$ where $a$ is the average degree of the graph and is usually small for real-world networks. 
The other key advantage of random walks over classic search-based sampling strategies is its time complexity. 
%
In particular, by imposing graph connectivity in the sample generation process, random walks provide a convenient mechanism to increase the effective sampling rate by reusing samples across different source nodes. By simulating a random walk of length $l>k$ we can generate $k$ samples for $l-k$ nodes at once due to the Markovian nature of the random walk. Hence, our effective complexity is $O\big(\frac{l}{k(l-k)}\big)$ per sample. For example, in Figure \ref{fig:bfsdfs} we sample a random walk $\{u, s_4, s_5, s_6, s_8, s_9\}$ of length $l=6$, which results in $N_S(u) =\{s_4, s_5, s_6\}$, $N_S(s_4)=\{s_5, s_6, s_8\}$ and $N_S(s_5)=\{s_6, s_8, s_9\}$.
Note that sample reuse can introduce some bias in the overall procedure. However, we observe that it greatly improves the efficiency.

%For example, in Figure \ref{fig:bfsdfs}, we see a random walk with $l=5$ in green starting from $u$. In order to satisfy the resource constraint $k$ for every source node using BFS and DFS, we need to run the sampling procedure $k|V|$ times, which gives a time complexity of $O(1)$ per sample. 
%While DFS is practically infeasible because of the high space and time preprocessing overhead, it is possible to exactly recover BFS in the random walk framework by simulating $k$ walks with $l=1$ from every source node. Doing so is similar to random walks with restarts, however, the runtime efficiency for walks with unit window size is $O(1)$. 


\subsubsection{The \nodevec algorithm}
 
\begin{algorithm}[h]
\begin{algorithmic}
\STATE \hspace{-0.15in}{\bf LearnFeatures}{ (Graph $G=(V, E, W)$, Dimensions $d$, Walks per node $r$, Walk length $l$, Context size $k$, Return $p$, In-out $q$)}
\STATE $\pi =$ PreprocessModifiedWeights($G, p, q$)
\STATE $G' = (V, E, \pi)$
\STATE Initialize $walks$ to Empty
\FOR {$iter = 1$ {\bfseries to} $r$} 
\FORALL {nodes $u \in V$}
\STATE $walk =$ node2vecWalk($G', u, l$)
\STATE Append $walk$ to $walks$
\ENDFOR
\ENDFOR
\STATE $f =$ StochasticGradientDescent($k$, $d$, $walks$)
\STATE {\bf return} $f$
\STATE \hspace{-0.15in} \hrulefill
\STATE \hspace{-0.15in}{\bf node2vecWalk}{ (Graph $G'= (V, E, \pi)$, Start node $u$, Length $l$)}
\STATE Inititalize $walk$ to $[u]$
\FOR {$walk\_iter = 1$ {\bfseries to} $l$}
\STATE $curr = walk[-1]$
\STATE $V_{curr} =$ GetNeighbors($curr$, $G'$)
\STATE $s =$ AliasSample($V_{curr}, \pi$)
\STATE Append $s$ to $walk$
\ENDFOR
\STATE {\bf return} $walk$
\end{algorithmic}
\caption{The \nodevec algorithm.}
\label{alg:node2vec}
\end{algorithm}

The pseudocode for \nodevec, is given in Algorithm~\ref{alg:node2vec}. In any random walk, there is an implicit bias due to the choice of the start node $u$. Since we learn representations for all nodes, we offset this bias by simulating $r$ random walks of fixed length $l$ starting from {\em every} node. At every step of the walk, sampling is done based on the transition probabilities $\pi_{vx}$.
% In our algorithm, every edge $(v,x)$ is associated with the following (unnormalized) transition probability.
% \begin{align*}
% 	\pi_{vx} =  \alpha_{pq}(t,x) \cdot w_{vx} 
% \end{align*}
% The individual terms in the above expression are defined as follows:
% \begin{itemize}[noitemsep,nolistsep]
% 	\item $w_{vx}$: Edge weight of edge $(v,x)$. %For unweighted graphs, $w_{vx}=1$ for all edges.
% 	\item $\alpha_{pq}(t,x)$: Search bias set using parameters $p$ and $q$.
% 	% \item $NI(x)$: Importance of node $x$. $NI(x) = \frac{PageRank(x)}{degree(x)}$.
% \end{itemize}
The transition probabilities $\pi_{vx}$ for the 2$^{\textrm{nd}}$ order Markov chain can be precomputed and hence, sampling of nodes while simulating the random walk can be done efficiently in $O(1)$ time using alias sampling. The three phases of \nodevec, \ie, preprocessing to compute transition probabilities, random walk simulations and optimization using SGD, are executed sequentially. Each phase is parallelizable and executed asynchronously, contributing to the overall scalability of \nodevec.

\nodevec is available at: \url{http://snap.stanford.edu/node2vec}. 

\begin{table}
\centering
\begin{tabular}{l|c|c}
\textbf{Operator} & \textbf{Symbol} & \textbf{Definition} \\ \hline
% Dot & $\centerdot$ & $ f(u) \centerdot f(v) = \sum_{i=1}^d f_i(u)f_i(v)$ \\
% L1 & $\| \cdot \|_1$ & $\|f(u) \cdot f(v)\|_1 = \sum_{i=1}^d |f_i(u)-f_i(v)|$ \\
% L2 & $\| \cdot \|_2$ & $\|f(u) \cdot f(v)\|_2 = \sum_{i=1}^d |f_i(u)-f_i(v)|^2$ \\ \hline
Average & $\boxplus$ & $[f(u) \boxplus f(v)]_i = \frac{f_i(u)+f_i(v)}{2}$ \\
Hadamard & $\boxdot$ & $[f(u) \boxdot f(v)]_i = f_i(u)*f_i(v)$ \\
Weighted-L1 & $\| \cdot \|_{\bar{1}}$ & $\|f(u) \cdot f(v)\|_{\bar{1}i} = |f_i(u)-f_i(v)|$ \\
Weighted-L2 & $\| \cdot \|_{\bar{2}}$ & $\|f(u) \cdot f(v)\|_{\bar{2}i} = |f_i(u)-f_i(v)|^2$
% Concat & $\odot$ &  $[f(u) \odot f(v)]_i = f_i(u)$ for $i\leq d$ \\
% & & $f_{i-d}(v)$ otherwise\\
\end{tabular}
\vspace{-0.2cm}
\caption{Choice of binary operators $\circ$ for learning edge features. The definitions correspond to the $i${th} component of $g(u,v)$.
%The first three are scalars, and the rest are vectors. All but the last are symmetric.
}
\vspace{-0.4cm}
\label{tab:edgeop}
\end{table} 

\subsection{Learning edge features}
The \nodevec algorithm provides a semi-supervised method to learn rich feature representations for nodes in a network. However, we are often interested in prediction tasks involving pairs of nodes instead of individual nodes. For instance, in link prediction, we predict whether a link exists between two nodes in a network. Since our random walks are naturally based on the connectivity structure between nodes in the underlying network, we extend them to pairs of nodes using a bootstrapping approach over the feature representations of the individual nodes.

Given two nodes $u$ and $v$, we define a binary operator $\circ$ over the corresponding feature vectors $f(u)$ and $f(v)$ in order to generate a representation $g(u,v)$ such that $g: V \times V \rightarrow \mathbb{R}^{d'}$ where $d'$ is the representation size for the pair $(u,v)$. We want our operators to be generally defined for any pair of nodes, even if an edge does not exist between the pair since doing so makes the representations useful for link prediction where our test set contains both true and false edges (\ie, do not exist). We consider several choices for the operator $\circ$ such that $d'= d$ which are summarized in Table \ref{tab:edgeop}.


